\chapter{Conclusion and Future Work}
\label{chapter:conclusion}

This final chapter synthesizes the architectural insights, empirical findings, and technical contributions established throughout this thesis. We present a formal summary of the research contributions (Section~\ref{sec:contributions}), explicitly answer the four research questions posed in Chapter~4 in light of the evidence gathered in Chapter~6 (Section~\ref{sec:rq-answers}), critically examine the operational limitations of the developed framework (Section~\ref{sec:limitations}), outline concrete directions for future research (Section~\ref{sec:future-work}), and close with concluding remarks (Section~\ref{sec:concluding-remarks}).

%==============================================================
\section{Summary of Technical Contributions}
\label{sec:contributions}
%==============================================================

This thesis has investigated the challenges of context drift and cross-turn dependency that characterize multi-turn dialogue in knowledge-intensive settings. Evaluated on the multi-domain MTRAGEval benchmark (SemEval-2026 Task~8), the proposed architecture indicates that retrieval quality at later dialogue turns is improved more effectively by targeted query diversity than by multi-retriever ensembles, and that deferring answer commitment to a multi-stage pipeline benefits grounded generation. The principal contributions are summarized below, each cross-referenced to its point of introduction and to the corresponding evidence.

\begin{itemize}
    \item \textbf{The Query-Diversity Paradigm} (Section~\ref{subsec:rewriting-engine}; evaluated in Section~\ref{sec:taskA-ablations}). Issuing five complementary, LLM-generated query reformulations over a single corpus-aligned sparse index (ELSER~v1) yields higher top-rank retrieval quality on the development set than ensembling heterogeneous retrievers; the complete configuration improves Recall@5 by $+25.7\%$ relative to the unaugmented baseline. Under the benchmark's official relevance judgments, adding alternative retrievers to the optimized configuration lowers Recall@5 (the retriever ensemble paradox, Section~\ref{sec:ensemble-paradox}), a result whose size may be affected by pooling bias (Section~\ref{sec:pooling-bias}).
    \item \textbf{Hierarchical Nested Rank Fusion (Nested RRF)} (Section~\ref{subsec:nested-rrf}). A two-level rank aggregation tree pre-aggregates the exploratory query tracks (HyDE, Chain-of-Thought, Anchor-Keyword) before merging them with the conservative ones (Minimal, Corpus-Specific). The complete retrieval configuration ranked first on the official Task~A leaderboard ($\text{nDCG@5}{=}0.5776$, $+20.5\%$ over the strongest baseline); the contribution of the hierarchical fusion itself, relative to a flat fusion of the same lists, was not isolated.
    \item \textbf{Span-Guided Generation} (Section~\ref{sec:generation-pipeline}; evaluated in Section~\ref{sec:taskB-ablations}). The generation pipeline replaces full-passage input with sentence-level verbatim evidence spans and selects between two drafted candidates using judge scores and an extractiveness shaping function $\phi(r_4)$ calibrated on the development references. Span extraction contributes the largest single gain in the generation ablation ($+0.034$ HM), and the pipeline ranked second on the official Task~B leaderboard ($\text{HM}{=}0.7698$).
    \item \textbf{A Multi-Layer Answerability Gate} (Section~\ref{sec:e2e-system}; evaluated in Section~\ref{sec:taskC-bottleneck}). Judges at the level of documents, evidence spans, and drafted answers are combined through confidence-weighted voting with dissenter override. On the development set the gate improves end-to-end HM over a single-judge gate ($+0.025$, Table~\ref{tab:taskC-ablation-en}), and qualitative inspection of its decisions points to the answer-level judge as the main source of the remaining acceptance bias.
    \item \textbf{An Empirical Characterization of the Answerability Bottleneck} (Section~\ref{sec:taskC-bottleneck}). Across all evaluated gate configurations, Unanswerable F1 stays below $25\%$ on the full development set (Table~\ref{tab:answerability-en}), with wide uncertainty given only $55$ unanswerable turns. On the test set, the gap between reference-grounded generation ($\text{HM}{=}0.7698$) and end-to-end RAG ($\text{HM}{=}0.5409$) points to answerability calibration under a tripled unanswerable rate as the decisive factor, while top-rank retrieval errors remain a substantial contributor: on the development set, where unanswerable turns are rare, they account for most of the inspected failures. Better-calibrated abstention is therefore a central open problem for multi-turn conversational RAG, alongside top-rank retrieval precision.
\end{itemize}

%==============================================================
\section{Answers to the Research Questions}
\label{sec:rq-answers}
%==============================================================

The four research questions formulated in Section~\ref{sec:rq} are answered here in light of the evidence of Chapter~6.

\paragraph{RQ1 --- Can query diversity substitute for retriever diversity?}
\textbf{Yes, under the benchmark's official judgments.} Query reformulation diversity over a single, corpus-aligned retriever outperforms retriever ensembling: the full five-strategy configuration reaches $\text{Recall@5}{=}0.607$ on the development set and $\text{nDCG@5}{=}0.5776$ on the official test set, while heterogeneous retriever ensembles improve deep recall (Recall@100) but lower top-rank precision once the base retriever is well optimized (Section~\ref{sec:ensemble-paradox}). Hypothesis~H1 is supported, with two qualifications: the comparison rests on development-set measurements, and it may be affected by pooling bias in MTRAG's relevance labels (Section~\ref{sec:pooling-bias}).

\paragraph{RQ2 --- Does evidence extraction suppress hallucination?}
\textbf{Yes, on the development set.} Verbatim evidence span extraction contributes the largest single improvement in the generation ablation ($+0.034$ HM), accompanied by a $+0.041$ gain in the reference-less faithfulness metric $RL_F$, as predicted by the faithfulness part of Hypothesis~H2: restricting the generator's input to a compressed, verifiable evidence set limits the space in which unsupported claims can be introduced. The high test-set $RL_F$ of the complete pipeline ($0.8971$) is consistent with this finding, although component-level effects cannot be measured on the test set. Fluency is not measured separately, but the reference-based metrics did not deteriorate on aggregate: a $+0.041$ change in one component of a three-way harmonic mean raises it by only a fraction of that amount, so the observed $+0.034$ HM gain implies that the reference-based components improved as well. Consequently, the asymmetry in favour of faithfulness that H2 also predicts cannot be confirmed, because per-metric changes of the reference-based metrics were not retained.

\paragraph{RQ3 --- How does distribution shift affect answerability gate sensitivity?}
\textbf{Substantially.} The precision--recall trade-off of answerability detection (Section~\ref{sec:pr-tradeoff}) is present in every gate configuration we evaluated, and none exceeds $25\%$ Unanswerable F1 on the full development set. The test set's unanswerable rate ($19.1\%$) is about three times the development rate ($6.5\%$) on which $\tau{=}0.7$ was calibrated, and the official Task~B$\rightarrow$C gap ($0.7698 \rightarrow 0.5409$) comes mostly from the drop in $RB_{alg}$ ($0.6327 \rightarrow 0.3998$). This makes answerability calibration the decisive end-to-end factor under the test distribution, although top-rank retrieval errors also contribute and the two effects cannot be separated without a per-category breakdown of the test scores. Hypothesis~H3 is partially supported: multi-layer judging improves over a single judge on the development set but does not eliminate acceptance bias, while the claim that distribution-adaptive calibration is required could not be tested directly with a single frozen submission.

\paragraph{RQ4 --- Does LLM-based reranking help once first-stage recall saturates?}
\textbf{Not meaningfully.} None of the three evaluated LLM-based reranking formulations (pointwise, listwise, generation-based) improved upon the weighted-RRF fusion of retriever and cross-encoder signals once first-stage $\text{Recall@100}$ exceeded $0.90$; as substitute ranking signals, all three lowered $\text{Recall@10}$ at a substantially higher API cost (Section~\ref{sec:reranking-role}), and used as a low-weight auxiliary signal they yield at most marginal Recall@5 gains (up to $+2$ points per domain) that do not justify this cost. Hypothesis~H4 is supported, with the practical implication that, in a mature retrieval pipeline, computational budget is better spent on query-side diversity than on LLM-based reranking.

%==============================================================
\section{System Limitations}
\label{sec:limitations}
%==============================================================

Despite its strong results on two of the three official leaderboards, the AILS-NTUA system and its evaluation carry several limitations that qualify the generality of the claims made in this thesis.

\begin{enumerate}
    \item \textbf{Development-set tuning and reporting.} All hyperparameters were tuned on the development set, and all ablations are reported on the same set; no held-out portion of the development data or conversation-level cross-validation was used, so the reported development numbers are likely optimistic. Each configuration was run once, without repetitions over random seeds, so small differences (of the order of $0.01$ HM or less) should not be over-interpreted, and the cumulative ablations depend on the order in which components were added.
    \item \textbf{Static decision thresholds under distribution shift.} The answerability confidence threshold ($\tau{=}0.7$), the extractiveness target band ($[0.28, 0.38]$), and the per-corpus Nested RRF weights (Table~\ref{tab:nested-rrf-weights-en}) were calibrated exclusively on the development set and were \emph{not} re-tuned for the test distribution during the frozen evaluation phase. This choice is a likely major driver of the Task~C gap, together with top-rank retrieval errors (Section~\ref{sec:taskC-bottleneck}). Whether a different threshold would have performed better on the test set could not be verified, since a single frozen configuration was submitted; on the development set, $\tau{=}0.6$ performs slightly worse than $\tau{=}0.7$ (Table~\ref{tab:threshold-sweep-en}).
    \item \textbf{Reference history in the evaluation.} MTRAGEval supplies the reference conversation history at every turn. Our evaluation therefore does not test robustness to the system's own earlier errors, which a deployed assistant would face (Section~\ref{sec:error-analysis}).
    \item \textbf{Dependence on proprietary, closed-source APIs.} The system relies exclusively on API-based models---GPT-4o, GPT-4o-mini, and DeepSeek-V3.2---for generation, judging, and query rewriting (Table~\ref{tab:routing-en}). This design introduces external runtime cost and variability outside the system's control and limits reproducibility, since neither the models' weights nor their training data are accessible. No component of the pipeline was fine-tuned on MTRAG data; every reported gain stems from prompt engineering and architectural design, so task-specific training could plausibly yield further improvements that this thesis does not explore.
    \item \textbf{Evaluation-side biases in the benchmark itself.} As noted in Section~6.2.1 and Section~\ref{sec:pooling-bias}, MTRAG's relevance annotations were constructed with ELSER-retrieved passages as the primary review pool, which may inflate ELSER's measured advantage over alternative retrievers and contribute to the ensemble paradox. Independently, the generation-quality metrics ($RB_{alg}$, $RB_{llm}$) depend on human reference answers and LLM judges that may reward specific phrasing styles over correct but differently worded responses---a limitation of reference-based generation evaluation in general.
    \item \textbf{Pipeline latency.} The sequential decomposition of the task into five retrieval tracks, cross-encoder reranking, span extraction, dual-candidate generation, and up to three answerability judges adds end-to-end latency, which we did not measure. Model routing (Section~\ref{sec:model-routing}) reduces the estimated monetary cost but not the sequential dependency chain, and the architecture would require additional orchestration engineering---parallelization of independent stages, caching across turns---before low-latency deployment.
    \item \textbf{Domain-balance mismatch between development and test conditions.} Beyond the unanswerable-rate shift, the test set under-represents FiQA ($15.2\%$ vs.\ $23.6\%$ at development time) and over-represents Govt ($31.0\%$ vs.\ $25.4\%$). Since the Nested RRF weights were tuned per domain on the development distribution, this compositional shift may compound the answerability threshold miscalibration discussed above.
\end{enumerate}

%==============================================================
\section{Avenues for Future Work}
\label{sec:future-work}
%==============================================================

The limitations enumerated above, together with the negative result of RQ4, point toward five concrete directions for future research.

\begin{itemize}
    \item \textbf{Distribution-Adaptive Threshold Calibration.} Rather than a fixed confidence threshold $\tau$, future systems could estimate the current unanswerable-rate regime online---for instance, via entropy tracking over the retrieved documents or by monitoring semantic divergence between successive turns---and adjust $\tau$ accordingly. Since answerability calibration is the decisive end-to-end factor under the test distribution, this is a high-leverage direction for future work; a per-category analysis of the official test scores would help quantify the expected gain.
    \item \textbf{Stronger Gate Baselines.} NLI-based sufficiency classifiers, learned classifiers over retrieval scores, and per-domain calibrated thresholds are natural baselines for the answerability gate that this thesis did not evaluate.
    \item \textbf{Robust Resolution of Cross-Turn Dependencies.} After current-turn retrieval errors, the largest category in our error analysis is unresolved dependency on earlier turns ($34\%$, Table~\ref{tab:error-decomp-en}). Better coreference and ellipsis resolution---for example, rewriting with explicit entity tracking across turns---targets this directly. In deployed settings, where the history contains the system's own outputs, mechanisms that audit and roll back earlier turns could additionally limit error cascades, a scenario that MTRAGEval does not evaluate.
    \item \textbf{Knowledge Distillation to Open-Weight Reasoning Models.} The structured reasoning traces produced by the multi-query rewriting engine and the multi-judge selection framework are a rich supervision signal that could be distilled into specialized, open-weight reasoning models (e.g., fine-tuned Llama or DeepSeek-R1 variants), compressing the multi-stage logic into a single model. This would address the closed-API dependency of Section~\ref{sec:limitations}, reducing inference cost and improving reproducibility.
    \item \textbf{Reinforcement Learning for Query Reformulation.} Rather than relying exclusively on prompt engineering, the query rewriter could be optimized via reinforcement learning from downstream retrieval signal---using $\text{nDCG@5}$ or Recall@5 directly as the reward (e.g., through PPO)---exploiting headroom on the query-formulation side, which the negative LLM-reranking result of RQ4 suggests is no longer available on the reranking side.
\end{itemize}

%==============================================================
\section{Concluding Remarks}
\label{sec:concluding-remarks}
%==============================================================

This thesis indicates that the success of multi-turn conversational RAG systems does not depend solely on model scaling. Shifting the architectural focus toward retrieval stability and decision calibration is a promising pathway toward reliable, verifiable information-seeking systems, supported by our development-set analyses and by the official rankings: treating conversational search as a query-diversity problem, and placing generation behind layered verification, the AILS-NTUA system ranked first in passage retrieval and second in reference-grounded generation, among $38$ and $26$ participants respectively.

At the same time, the eleventh-place result on the end-to-end task---only marginally above the strongest organizer baseline---is informative in its own right. Our analysis identifies \emph{answerability calibration} as the decisive open problem between current multi-turn RAG systems and reliable deployment, with a weight that grows with the share of unanswerable requests, while top-rank retrieval precision remains a substantial secondary constraint. The central hypothesis of this thesis (Section~\ref{sec:rq}) is thus supported in its retrieval and answerability parts, while its claim about model capacity was not tested directly: retrieval stability can be improved through disciplined query diversity without model scaling, but end-to-end reliability remains bounded by how well the decision to answer or to abstain is calibrated to a possibly shifting distribution of requests.
